\section{Introduction}\label{sec:intro}

A prognostic model creates value through the maintenance actions it triggers. For a component that is replaced once, near the end of its life, the action is usually triggered by an \emph{end-of-life alarm}: a policy watches the unit cycle by cycle and raises one alarm, which schedules a replacement after a lead time for parts, crew and access. The first reliability requirement on such an alarm is therefore not the accuracy of a remaining useful life (RUL) estimate at some cycle chosen by the evaluator, but the probability that the unit fails before the replacement the alarm scheduled. This paper gives that probability a distribution-free, finite-sample guarantee.

The usual ways of evaluating and calibrating prognostics do not control it. Accuracy metrics and asymmetric acceptance windows score an estimate at a time, and a window score at the alarm mixes the estimate with the alarm time~\cite{saxena2008,saxena2010}. Decision-oriented metrics, such as the metric $M$ of Kamariotis et al.~\cite{kamariotis2024}, score and tune policies by long-run cost but give no guarantee for a new unit. Conformal prediction gives RUL intervals with distribution-free coverage~\cite{javanmardi2023,robinson2026,wang2025robustuq}, but the coverage holds at a cycle chosen independently of the data, whereas an alarm acts at the first cycle at which the estimate, or its lower bound, crosses a level: a stopping time at which the estimate is, by selection, at its most optimistic. Two further obstacles stand between any calibration and the field. Fleet data are censored, because most units are replaced before they fail; and populations change, from one production batch or operating campaign to the next.

This paper calibrates the alarm itself. For a threshold trigger on any causal signal, the notice left by the alarm is monotone in the warning level, so each calibration unit has a single \emph{critical level}: the lowest warning level at which its alarm would still have left more than the lead time. The split-conformal quantile of these critical levels is a warning level with a finite-sample, distribution-free guarantee that a new exchangeable unit is replaced before it fails with probability at least $1-\alpha$. The contributions are:

\begin{enumerate}[leftmargin=1.4em,itemsep=2pt]
\item \textbf{Theory.} A closed-form critical level for debounced threshold triggers (Lemma~\ref{lem:crit}); the notice guarantee with a matching upper bound (Theorem~\ref{thm:main}) and its training-conditional form, which says how far the failure probability of one calibrated alarm can stray from $\alpha$ (Corollary~\ref{cor:pac}); a decision-aware choice of the level that keeps the guarantee (Corollary~\ref{cor:da}); a cost bound and a certification floor that state what $n$ calibration units can certify (Proposition~\ref{prop:cost}); a pathwise censoring bound under which the guarantee holds for \emph{any} censoring mechanism, so that units replaced before failure can be used for calibration (Lemma~\ref{lem:cens}, Theorem~\ref{thm:cens}); a cross-fitted variant that calibrates on all training units (Theorem~\ref{thm:cvplus}); and an online update that bounds the long-run failure rate for any sequence of units, without exchangeability (Theorem~\ref{thm:online}). A supporting result (Proposition~\ref{prop:window}) explains why window scores of threshold alarms measure timing and cannot certify notice.
\item \textbf{Method.} The conformal notice alarm (CNA) and its decision-aware, cross-fitted and online variants (Algorithm~\ref{alg:cna}), which wrap any signal, including a learned regressor and a physically motivated capacity trend.
\item \textbf{Evidence.} A protocol and four amendments, each hashed before its experiments; seven datasets (PHM08, C-MAPSS FD001--FD004, N-CMAPSS and 123 lithium-ion cells in three batches); comparisons with interval-calibrated alarms, window- and cost-tuned alarms, the probability-threshold policy of Kamariotis et al.\ (also in their own decision setting and metric), age replacement and a capacity-trend baseline; four learners and 27 alarm settings; three censoring mechanisms; and transfer to held-out battery batches and N-CMAPSS subsets.
\end{enumerate}

% stub sec1_findings


Section~\ref{sec:related} reviews related work, Section~\ref{sec:problem} states the problem, Section~\ref{sec:theory} gives the theory and Section~\ref{sec:method} the method. Section~\ref{sec:design} describes the design, Section~\ref{sec:results} the results and Section~\ref{sec:discussion} their implications and limits.
